% 字体、页面、章标题、页眉页脚、目录和图表样式。
\usepackage{geometry,fancyhdr}
\usepackage[indentafter]{titlesec}
\usepackage{titletoc}
\usepackage{graphicx,multicol,tabularx,booktabs,longtable}
\usepackage{amsmath,amssymb,bm}
\usepackage{caption}
\usepackage[numbers,square,comma,super,sort&compress]{natbib}
\usepackage{etoolbox}
\usepackage[unicode,hidelinks]{hyperref}

\setmainfont[Path=fonts/,BoldFont=timesbd.ttf,ItalicFont=timesi.ttf,BoldItalicFont=timesbi.ttf]{times.ttf}
\setCJKmainfont[Path=fonts/,AutoFakeBold=2]{song.ttf}
\setCJKsansfont[Path=fonts/]{hei.ttf}
\setCJKmonofont[Path=fonts/]{fangsong.ttf}
\setCJKfamilyfont{zhsong}[Path=fonts/,AutoFakeBold=2]{song.ttf}
\setCJKfamilyfont{zhhei}[Path=fonts/]{hei.ttf}
\setCJKfamilyfont{zhkai}[Path=fonts/]{kai.ttf}
\setCJKfamilyfont{zhfs}[Path=fonts/]{fangsong.ttf}
% 校名为中文，必须使用 xeCJK 字体命令。
\newCJKfontfamily\jianqifont[Path=fonts/]{jianqi.ttf}
\providecommand\songti{\CJKfamily{zhsong}}
\providecommand\heiti{\CJKfamily{zhhei}}
\providecommand\kaishu{\CJKfamily{zhkai}}
\providecommand\fangsong{\CJKfamily{zhfs}}

\xeCJKDeclareCharClass{CJK}{"2573}

% 摘要标题采用Word实际使用的SimHei，两个空格也使用同一字体。
\newfontfamily\wordheiLatin[Path=fonts/]{simhei.ttf}
\setCJKfamilyfont{wordhei}[Path=fonts/]{simhei.ttf}

% A4双面；连续换页，不插入奇偶页补白。
\geometry{asymmetric,top=3.3cm,bottom=3.3cm,left=3cm,right=2.8cm,
  headheight=15.6bp,headsep=\dimexpr .7cm-15.6bp\relax,footskip=.7cm}
\let\cleardoublepage\clearpage
\linespread{1}
\makeatletter
% 这一处底层命令用于固定正文为五号、20磅行距。
\renewcommand{\normalsize}{\@setfontsize\normalsize{10.5bp}{20bp}}
\makeatother
\normalsize
\setlength{\parindent}{21bp}
\setlength{\parskip}{0pt}
\setcounter{tocdepth}{2}
\setcounter{secnumdepth}{3}
\ctexset{chapter={lofskip=0pt,lotskip=0pt}}
\titleformat{\chapter}{\centering\wordheiLatin\CJKfamily{wordhei}\fontsize{15bp}{20bp}\selectfont}
  {第\chinese{chapter}章}{1em}{}
\titleformat*{\section}{\heiti\fontsize{14bp}{20bp}\selectfont}
\titleformat*{\subsection}{\heiti\fontsize{12bp}{20bp}\selectfont}
\titleformat*{\subsubsection}{\heiti\fontsize{12bp}{20bp}\selectfont}
% Word章标题：段前后各1.5行、固定20磅；补偿titlesec页首行框后对齐实际位置。
\titlespacing*{\chapter}{0pt}{9.454bp}{23.44bp}
\titlespacing*{\section}{0pt}{7.8bp}{7.8bp}
\titlespacing*{\subsection}{0pt}{7.8bp}{7.8bp}
\titlespacing*{\subsubsection}{0pt}{7.8bp}{7.8bp}

% 页眉与页码：摘要起启用；奇数页学校，偶数页论文题目。
% 页码按Word：奇数页右下角（OR），偶数页左下角（EL）。
\fancypagestyle{thesis}{
  \fancyhf{}
  \fancyhead[C]{\songti\fontsize{10.5bp}{20bp}\selectfont
    % 保留页眉行框，仅上移文字基线；避免调整横线后压住文字。
    \strut\smash{\raisebox{3.461bp}{%
      \ifodd\value{page}南京航空航天大学专业学位硕士学位论文\else\ThesisTitle\fi}}}

  \fancyfoot[OR,EL]{\fontsize{7.5bp}{10bp}\selectfont\thepage}
  % Word页眉：单倍行距、段前后0；下边框0.75磅，距文字1磅。
  % fancyhdr默认横线偏低；上移5.381磅，保持摘要标题/正文位置不动。
  \renewcommand{\headrulewidth}{.75bp}
  \renewcommand{\headrule}{%
    \vskip-5.381bp
    \hrule height .75bp width\headwidth
    \vskip4.631bp
  }
  \renewcommand{\footrulewidth}{0pt}
}
\makeatletter
\let\ps@plain\ps@thesis
\makeatother
\assignpagestyle{\chapter}{thesis}

% 前置页共用标题：与已核对的目录页保持同一位置和上下间距。
\newcommand{\FrontHeading}[1]{%
  \begingroup
    \xeCJKsetup{CJKecglue=0pt}
    \titleformat{\chapter}{\centering\wordheiLatin\CJKfamily{wordhei}\fontsize{15bp}{20bp}\selectfont}{}{0pt}{}
    \titlespacing*{\chapter}{0pt}{-8.786bp}{3.684bp}
    \chapter*{#1}
  \endgroup
}

% 自动目录：按Word的标题、层级缩进与文档网格设置。
\renewcommand{\contentsname}{目\hbox{\wordheiLatin\char32\char32}录}
\newcommand{\TocFont}{\songti\fontsize{10.5bp}{15.6bp}\selectfont}
\contentsmargin{5.395bp} % 页码右边缘按Word制表位对齐（约521磅）
\newcommand{\TocPage}{\hbox to5.114bp{\leaders\hbox{.}\hfill}\contentspage\hspace*{-5.114bp}}
\titlecontents{chapter}[56.35bp]{\addvspace{1.92bp}\TocFont}
  {\contentslabel{36.75bp}}{\hspace*{-36.75bp}}
  {\titlerule*[2.625bp]{.}\TocPage}[\vspace{1.92bp}]
\titlecontents{section}[52.5bp]{\TocFont}
  {\contentslabel[\wordheiLatin\thecontentslabel]{32.9bp}}{}
  {\titlerule*[2.625bp]{.}\TocPage}
\titlecontents{subsection}[70.87bp]{\TocFont}
  {\contentslabel[\wordheiLatin\thecontentslabel]{28.87bp}}{}
  {\titlerule*[2.625bp]{.}\TocPage}
\DeclareCaptionFont{thesis}{\songti\fontsize{10.5bp}{20bp}\selectfont}
\captionsetup{font=thesis,labelsep=quad,justification=centering}
\captionsetup[table]{position=top}
\newcommand{\figtabnote}[1]{\par{\kaishu\fontsize{10.5bp}{20bp}\selectfont（注：#1）}\par}
\graphicspath{{\NUAAContentDirectory/fig/}{fig/}}
\hypersetup{pdftitle={\ThesisTitle},pdfauthor={\AuthorName}}

\newcommand{\ThesisAppendixHeading}[1]{%
  \begingroup
  \titlespacing*{\chapter}{0pt}{17bp}{16.5bp}%
  \titleformat{\chapter}{\centering\heiti\fontsize{15bp}{26bp}\selectfont}{}{0pt}{}%
  \chapter*{附录 #1}%
  \addcontentsline{toc}{chapter}{附录 #1}%
  \endgroup
}
